% !PRINTED 704
\weberpage{704}
From this it follows that in the product (6) the partial products
\[
\prod \frac{1}{1-p^{-sf}},
\]
which extend over all primes $p$ for which $f>1$, still remain finite
for $s=1$ as well. Since on the other hand $g, h$ have definite
values different from zero, there follows from (7) the important
theorem:

I. \gesperrt{If $\mathfrak{p}$ runs through the totality of prime ideals
of the first degree of any algebraic field $\Omega$, then the product}
\weq{10}{(s-1)\prod\frac{1}{1-N(\mathfrak{p})^{-s}}}
\gesperrt{has for $s=1$ a finite limiting value different from
zero.}

This property is retained even when, in the formation of the
product, an arbitrary \gesperrt{finite} number of prime ideals of the first
degree is omitted.

The norms $N(\mathfrak{p})$ in the formula (10) are equal to natural
primes $p$. But a prime $p$ occurs in it as often as it contains
prime factors of the first degree in $\Omega$, hence at most
$n$ times. If $\Omega$ is a \gesperrt{normal field}, then every prime $p$ that
is decomposable into prime factors of the first degree does in
fact occur in it exactly $n$ times, and for the product (10) we
may also put:
\[
(s-1)\prod\frac{1}{(1-p^{-s})^n},
\]
in which the product $\prod$ extends over all primes $p$ that are
\gesperrt{decomposable into prime factors of the first degree} (\S.\,160).
% [source note] The assertion of exactly n distinct degree-one prime factors requires an unramified prime. For Q(i), 2 ramifies as a unit times (1+i)^2 and has only one distinct degree-one prime. Omitting the finitely many ramified Euler factors preserves the finite nonzero limiting-value conclusion, but the displayed products are not literally identical without that qualification.

An immediate consequence of Theorem I. is this:

\gesperrt{In every algebraic field there exist infinitely
many prime ideals of the first degree.}

\begin{center}
\rule{6em}{0.4pt}
\end{center}
